\subsection{KMP: buscar un patrón}\label{rec:kmp}

Encuentra todas las apariciones de un patrón, incluso solapadas. Entrada:
\codigo{texto patron}. Salida: posiciones 0-based.

\tablafunc{
\codigo{pref(s)} & \codigo{pi[i]}: borde más largo de \codigo{s[0..i]} & $O(n)$ \\
\codigo{kmp(text, pat)} & posiciones donde aparece \codigo{pat} & $O(n+m)$ \\
}

\copiar{el bloque \texttt{COPIA KMP}. Devuelve las posiciones 0-based.}
\lstinputlisting{codigo/kmp.cpp}

El valor \codigo{pi[m-1]} es el borde más largo del patrón; el período más
pequeño es \codigo{m - pi[m-1]}. Si ese período divide a \codigo{m}, el
patrón es repetición de un bloque.